% GENERATED FILE. Edit canonical lessons, then run npm run generate:books.
% canonical-source-hash: fnv1a64:11ad624dd535afed
% canonical-lessons: UR-C265-us-ne-khaya, UR-C265-main-jaun-ga, UR-W265-waw-hamza, UR-R265-who-did-it

\chapter{Who Did It, and What Will Happen}
\label{ch:ur-who-did-it-and-what-will-happen}
\hlchaptermodality{265}

\begin{chapteropening}
\textbf{By the end of this chapter:} \emph{I can say who did something, with ne and a verb that follows the thing (us ne roṭī khāī), and what will happen, for every person (karūṅ gā, kare gā, kareṅ ge).}

Complete UR-R265-who-did-it: say who did something, and what will happen.
\end{chapteropening}

\section[\texorpdfstring{\ur{اس} \ur{نے} \ur{کھایا} (us ne khāyā)}{اس نے کھایا (us ne khāyā)}]{\ur{اس} \ur{نے} \ur{کھایا} (us ne khāyā) — who did it}

\label{lesson:UR-C265-us-ne-khaya}



\begin{quote}
Say \textquotedblleft{}he went home\textquotedblright{}. (\emph{voh ghar gayā}.) Going has no object. Eating has one: \textquotedblleft{}he ate bread\textquotedblright{}. That past is built another way.
\end{quote}

\subsection*{You'll want to know: \ur{نے} after the doer}
\begin{quote}
\textbf{\ur{میں} \ur{نے} \ur{کھایا۔}}
\end{quote}

\begin{quote}
\emph{maiṅ ne khāyā.} — \textquotedblleft{}I ate.\textquotedblright{}
\end{quote}

With a verb that has an object, the doer takes \textbf{\ur{نے}} (\emph{ne}), and the verb follows the THING, not the doer:

\begin{quote}
\textbf{\ur{اس} \ur{نے} \ur{روٹی} \ur{کھائی۔}}
\end{quote}

\begin{quote}
\emph{us ne roṭī khāī.} — \textquotedblleft{}He ate bread.\textquotedblright{}
\end{quote}

\textbf{\ur{روٹی}} (\emph{roṭī}) is feminine, so \textbf{\ur{کھائی}} (\emph{khāī}), even for a man. \textbf{\ur{وہ}} (\emph{voh}) becomes \textbf{\ur{اس}} (\emph{us}) before \textbf{\ur{نے}}:

\begin{quote}
\textbf{\ur{سارا} \ur{نے} \ur{خط} \ur{لکھا۔}}
\end{quote}

\begin{quote}
\emph{Sārā ne khat likhā.} — \textquotedblleft{}Sara wrote a letter.\textquotedblright{}
\end{quote}

\begin{quote}
\textbf{\ur{سارا} \ur{نے} \ur{کتاب} \ur{پڑھی۔}}
\end{quote}

\begin{quote}
\emph{Sārā ne kitāb paṛhī.} — \textquotedblleft{}Sara read a book.\textquotedblright{}
\end{quote}

\textbf{\ur{خط}} (\emph{khat}) is masculine: \textbf{\ur{لکھا}}. \textbf{\ur{کتاب}} (\emph{kitāb}) is feminine: \textbf{\ur{پڑھی}}.

\noindent
\begin{tabularx}{\linewidth}{@{}>{\raggedright\arraybackslash}X>{\raggedright\arraybackslash}X@{}}
\toprule
\textbf{the doer} & \textbf{with \ur{نے}} \\
\midrule
\textbf{\ur{میں}} & \textbf{\ur{میں} \ur{نے}} (\emph{maiṅ ne}) \\
\textbf{\ur{ہم}}, \textbf{\ur{آپ}}, \textbf{\ur{تم}} & \textbf{\ur{ہم} \ur{نے}}, \textbf{\ur{آپ} \ur{نے}}, \textbf{\ur{تم} \ur{نے}} \\
\textbf{\ur{وہ}}, one & \textbf{\ur{اس} \ur{نے}} (\emph{us ne}) \\
\textbf{\ur{وہ}}, several & \textbf{\ur{انہوں} \ur{نے}} (\emph{unhoṅ ne}) \\
\bottomrule
\end{tabularx}

\subsection*{You'll want to know: \ur{کیا}, \ur{دیا}, \ur{لیا}, \ur{کہا}, and \ur{مجھے}}
Four common pasts: \textbf{\ur{کیا}} (\emph{kiyā}), \textquotedblleft{}did\textquotedblright{}, from \textbf{\ur{کرنا}}; \textbf{\ur{دیا}} (\emph{diyā}), \textquotedblleft{}gave\textquotedblright{}; \textbf{\ur{لیا}} (\emph{liyā}), \textquotedblleft{}took\textquotedblright{}; \textbf{\ur{کہا}} (\emph{kahā}), \textquotedblleft{}said\textquotedblright{}. \textbf{\ur{کیا}} \emph{kiyā} is spelled like \textbf{\ur{کیا}} \emph{kyā}, \textquotedblleft{}what\textquotedblright{}: the romanization tells them apart.

\begin{quote}
\textbf{\ur{میں} \ur{نے} \ur{کام} \ur{کیا۔}}
\end{quote}

\begin{quote}
\emph{maiṅ ne kām kiyā.} — \textquotedblleft{}I did the work.\textquotedblright{}
\end{quote}

When the thing is a person with \textbf{\ur{کو}} (\emph{ko}), the verb stays masculine singular:

\begin{quote}
\textbf{\ur{اس} \ur{نے} \ur{سارا} \ur{کو} \ur{دیکھا۔}}
\end{quote}

\begin{quote}
\emph{us ne Sārā ko dekhā.} — \textquotedblleft{}He saw Sara.\textquotedblright{}
\end{quote}

\textbf{\ur{مجھے}} (\emph{mujhe}) is \textbf{\ur{میں}} with \textbf{\ur{کو}} inside it, \textquotedblleft{}me, to me\textquotedblright{}; so too \textbf{\ur{اسے}} (\emph{use}), \textbf{\ur{ہمیں}} (\emph{hameṅ}), \textbf{\ur{تمہیں}} (\emph{tumheṅ}), \textbf{\ur{انہیں}} (\emph{unheṅ}):

\begin{quote}
\textbf{\ur{سارا} \ur{نے} \ur{مجھے} \ur{دیکھا۔}}
\end{quote}

\begin{quote}
\emph{Sārā ne mujhe dekhā.} — \textquotedblleft{}Sara saw me.\textquotedblright{}
\end{quote}

\begin{quote}
\textbf{\ur{اس} \ur{نے} \ur{مجھے} \ur{پانی} \ur{دیا۔}}
\end{quote}

\begin{quote}
\emph{us ne mujhe pānī diyā.} — \textquotedblleft{}He gave me water.\textquotedblright{}
\end{quote}

One verb with an object goes the other way: \textbf{\ur{لانا}} (\emph{lānā}), \textquotedblleft{}to bring\textquotedblright{}, takes no \textbf{\ur{نے}}, like \textbf{\ur{آنا}} (\emph{ānā}):

\begin{quote}
\textbf{\ur{وہ} \ur{پانی} \ur{لایا۔}}
\end{quote}

\begin{quote}
\emph{voh pānī lāyā.} — \textquotedblleft{}He brought water.\textquotedblright{}
\end{quote}

\subsection*{Your turn}
Say these aloud:

\begin{itemize}
\raggedright
  \item \textquotedblleft{}I drank tea\textquotedblright{} — \emph{maiṅ ne chāy pī}
  \item \textquotedblleft{}Sara drank tea\textquotedblright{} — \emph{Sārā ne chāy pī}
  \item \textquotedblleft{}Sara drank water\textquotedblright{} — \emph{Sārā ne pānī piyā}
\end{itemize}

\begin{tcolorbox}[breakable,colback=teal!4,colframe=teal!35!black,title={Before you move on}]
\textquotedblleft{}He ate bread\textquotedblright{}? (\textbf{\ur{اس} \ur{نے} \ur{روٹی} \ur{کھائی}}.) Why \textbf{\ur{کھائی}}? (\textbf{\ur{روٹی}} is feminine.) \textquotedblleft{}She saw me\textquotedblright{}? (\textbf{\ur{اس} \ur{نے} \ur{مجھے} \ur{دیکھا}}.) And \textbf{\ur{وہ} \ur{گیا}}: does it take \textbf{\ur{نے}}? (No: going has no object.)
\end{tcolorbox}
\section[\texorpdfstring{\ur{میں} \ur{جاؤں} \ur{گا} (maiṅ jāūṅ gā)}{میں جاؤں گا (maiṅ jāūṅ gā)}]{\ur{میں} \ur{جاؤں} \ur{گا} (maiṅ jāūṅ gā) — what will happen}

\label{lesson:UR-C265-main-jaun-ga}



\begin{quote}
You met one future long ago: \textbf{\ur{میں} \ur{پڑھوں} \ur{گا}} (\emph{maiṅ paṛhūṅ gā}), \textquotedblleft{}I will read\textquotedblright{}. Say it, then a woman's \textbf{\ur{پڑھوں} \ur{گی}} (\emph{paṛhūṅ gī}). This lesson gives every person.
\end{quote}

\subsection*{You'll want to know: the future, person by person}
\begin{quote}
\textbf{\ur{میں} \ur{جاؤں} \ur{گا۔}}
\end{quote}

\begin{quote}
\emph{maiṅ jāūṅ gā.} — \textquotedblleft{}I will go.\textquotedblright{}
\end{quote}

\begin{quote}
\textbf{\ur{وہ} \ur{جائے} \ur{گی۔}}
\end{quote}

\begin{quote}
\emph{voh jāe gī.} — \textquotedblleft{}She will go.\textquotedblright{}
\end{quote}

A short ending follows the person, and \textbf{\ur{گا}} (\emph{gā}), \textbf{\ur{گی}} (\emph{gī}) or \textbf{\ur{گے}} (\emph{ge}) follows the man, the woman or the group. The book writes it apart, as it wrote \textbf{\ur{پڑھوں} \ur{گا}}:

\begin{quote}
\textbf{\ur{ہم} \ur{کل} \ur{جائیں} \ur{گے۔}}
\end{quote}

\begin{quote}
\emph{ham kal jāeṅ ge.} — \textquotedblleft{}We will go tomorrow.\textquotedblright{}
\end{quote}

\begin{quote}
\textbf{\ur{میں} \ur{کام} \ur{کروں} \ur{گی۔}}
\end{quote}

\begin{quote}
\emph{maiṅ kām karūṅ gī.} — \textquotedblleft{}I will do the work.\textquotedblright{}
\end{quote}

\begin{quote}
\textbf{\ur{تم} \ur{کب} \ur{آؤ} \ur{گے؟}}
\end{quote}

\begin{quote}
\emph{tum kab āo ge?} — \textquotedblleft{}When will you come?\textquotedblright{}
\end{quote}

\noindent
\begin{tabularx}{\linewidth}{@{}>{\raggedright\arraybackslash}X>{\raggedright\arraybackslash}X>{\raggedright\arraybackslash}X@{}}
\toprule
\textbf{who} & \textbf{\ur{کرنا}} & \textbf{\ur{جانا}} \\
\midrule
\textbf{\ur{میں}} & \textbf{\ur{کروں} \ur{گا}} (\emph{karūṅ gā}) & \textbf{\ur{جاؤں} \ur{گا}} (\emph{jāūṅ gā}) \\
\textbf{\ur{تم}} & \textbf{\ur{کرو} \ur{گے}} (\emph{karo ge}) & \textbf{\ur{جاؤ} \ur{گے}} (\emph{jāo ge}) \\
\textbf{\ur{وہ}}, one & \textbf{\ur{کرے} \ur{گا}} (\emph{kare gā}) & \textbf{\ur{جائے} \ur{گا}} (\emph{jāe gā}) \\
\textbf{\ur{ہم}}, \textbf{\ur{آپ}}, \textbf{\ur{وہ}} & \textbf{\ur{کریں} \ur{گے}} (\emph{kareṅ ge}) & \textbf{\ur{جائیں} \ur{گے}} (\emph{jāeṅ ge}) \\
\bottomrule
\end{tabularx}

\textquotedblleft{}Won't\textquotedblright{} puts \textbf{\ur{نہیں}} before it: \textbf{\ur{میں} \ur{نہیں} \ur{جاؤں} \ur{گا}} (\emph{maiṅ nahīṅ jāūṅ gā}).

\subsection*{You'll want to know: \ur{ہو} \ur{گا}, \textquotedblleft{}will be\textquotedblright{}}
\begin{quote}
\textbf{\ur{کل} \ur{بارش} \ur{ہو} \ur{گی۔}}
\end{quote}

\begin{quote}
\emph{kal bārish ho gī.} — \textquotedblleft{}It will rain tomorrow.\textquotedblright{}
\end{quote}

\textbf{\ur{بارش}} (\emph{bārish}) is feminine: \textbf{\ur{ہو} \ur{گی}}. For a man, \textbf{\ur{ہو} \ur{گا}}; for a group, \textbf{\ur{ہوں} \ur{گے}} (\emph{hoṅ ge}):

\begin{quote}
\textbf{\ur{وہ} \ur{شاید} \ur{گھر} \ur{پر} \ur{ہو} \ur{گا۔}}
\end{quote}

\begin{quote}
\emph{voh shāyad ghar par ho gā.} — \textquotedblleft{}He will probably be at home.\textquotedblright{}
\end{quote}

\subsection*{Your turn}
Say these aloud:

\begin{itemize}
\raggedright
  \item \textquotedblleft{}I will come tomorrow\textquotedblright{} — \emph{maiṅ kal āūṅ gā}
  \item \textquotedblleft{}She will come tomorrow\textquotedblright{} — \emph{voh kal āe gī}
  \item \textquotedblleft{}They will come tomorrow\textquotedblright{} — \emph{voh kal āeṅ ge}
\end{itemize}

\begin{tcolorbox}[breakable,colback=teal!4,colframe=teal!35!black,title={Before you move on}]
\textquotedblleft{}I will do it\textquotedblright{}? (\textbf{\ur{میں} \ur{کروں} \ur{گا}}, or \textbf{\ur{کروں} \ur{گی}}.) \textquotedblleft{}You will go\textquotedblright{}, to a friend? (\textbf{\ur{تم} \ur{جاؤ} \ur{گے}}.) \textquotedblleft{}He will be at home\textquotedblright{}? (\textbf{\ur{وہ} \ur{گھر} \ur{پر} \ur{ہو} \ur{گا}}.)
\end{tcolorbox}
\section[\texorpdfstring{\ur{ؤ} (vāʾo hamza)}{ؤ (vāʾo hamza)}]{\ur{ؤ} (vāʾo hamza) — vāʾo wearing the hamza, inside \ur{جاؤں}}

\label{lesson:UR-W265-waw-hamza}



\begin{quote}
Before the new one: write \textbf{\ur{و}} and \textbf{\ur{ئ}} once each, and name them. (\emph{vāʾo}, \emph{hamza}.)

Now say \textbf{\ur{میں} \ur{جاؤں} \ur{گا}} (\emph{maiṅ jāūṅ gā}), \textquotedblleft{}I will go\textquotedblright{}. You have read \textbf{\ur{جاؤں}} (\emph{jāūṅ}) in every person's future. This lesson takes one letter out of it and puts it in your hand.
\end{quote}

\subsection*{Script you'll notice: \ur{ؤ}}
\textbf{\ur{ؤ}} — \emph{vāʾo hamza}.

It is inside \textbf{\ur{جاؤں}} (\emph{jāūṅ}) and \textbf{\ur{آؤ}} (\emph{āo}). Say \emph{jā-ūṅ}: two vowels side by side. \textbf{\ur{ؤ}} is \textbf{\ur{و}} wearing the same small \textbf{hamza} you wrote on \textbf{\ur{ئ}}: there it sat on a tooth, here it sits on \textbf{\ur{و}}. It marks the place where one vowel ends and the next begins. Like \textbf{\ur{و}}, it does not join the letter after it.

\subsection*{Writing: \ur{ؤ}}
Put your pen on \ur{ؤ} and follow its line. Copy the shape you can see — slowly, and larger than it is printed.

\begin{quote}
This book does not yet tell you \textbf{where to start this letter or which way to travel}. It is not written down here until it can be written down with a source. Copying what is in front of you needs no such source, and it is how the shape gets into your hand in the meantime.
\end{quote}

\subsection*{Your turn}
\begin{itemize}
\raggedright
  \item \emph{Look:} at \textbf{\ur{جاؤں}} and point at \ur{ؤ} inside it
  \item \emph{Trace it:} \ur{ؤ} three times, saying \emph{jā-ūṅ} as you finish each one
  \item \emph{Write it:} \textbf{\ur{جاؤں}} once, joined, with \ur{ؤ} in its place
  \item \emph{Read it:} \textbf{\ur{جاؤں}} aloud once more
\end{itemize}

\begin{tcolorbox}[breakable,colback=teal!4,colframe=teal!35!black,title={Before you move on}]
Which letter is this — \ur{ؤ}? (\textbf{vāʾo hamza}, \emph{jā-ūṅ}.) Which word did it come from? (\textbf{\ur{جاؤں}}, \emph{jāūṅ}, in \textbf{\ur{میں} \ur{جاؤں} \ur{گا}}.)
\end{tcolorbox}
\section[(review)]{Review — who did it, and what will happen}

\label{lesson:UR-R265-who-did-it}



\begin{quote}
Say \textquotedblleft{}she went home\textquotedblright{} and \textquotedblleft{}I used to go to school\textquotedblright{}. (\emph{voh ghar gaī}, \emph{maiṅ iskūl jātā thā}.)
\end{quote}

\subsection*{Your turn}
\begin{itemize}
\raggedright
  \item \emph{Say it:} He ate bread. Sara wrote a letter.
  \item \emph{Say it:} I did the work, and she gave me tea.
  \item \emph{Say it:} Sara saw me at the market.
  \item \emph{Say it:} Tomorrow I will go home. We will go too.
  \item \emph{Say it:} It will rain. He will probably be at home.
  \item \emph{Write it:} \textbf{\ur{جاؤں}} and \textbf{\ur{آؤ}}, then read them aloud
\end{itemize}

\begin{tcolorbox}[breakable,colback=teal!4,colframe=teal!35!black,title={Before you move on}]
In \textbf{\ur{سارا} \ur{نے} \ur{خط} \ur{لکھا}}, which word sets \textbf{\ur{لکھا}}? (\textbf{\ur{خط}}.) Put \textbf{\ur{کتاب}} in place of \textbf{\ur{خط}}: does the verb change? (Yes: \textbf{\ur{سارا} \ur{نے} \ur{کتاب} \ur{لکھی}}.) Now the other way round, the future: \textquotedblleft{}she will go\textquotedblright{}? (\textbf{\ur{وہ} \ur{جائے} \ur{گی}}.) \textquotedblleft{}We will go\textquotedblright{}? (\textbf{\ur{ہم} \ur{جائیں} \ur{گے}}.)
\end{tcolorbox}
